\section{Discussion}
\label{sec:discussion}

This section explains the results (Section~\ref{subsec:why}), draws a recommendation for tuning studies (Section~\ref{subsec:implications}), considers deployment (Section~\ref{subsec:deployment}) and states the limits of the evidence (Section~\ref{subsec:limitations}).

\subsection{Why PSO Did Not Beat Random Search}
\label{subsec:why}

In \genone, the defaults left little room for any search. On the proxy they scored 0.4765 \mapfifty, above 28 of the 32 random-search evaluations and 21 of the 32 PSO evaluations (Section~\ref{subsec:res-search}), and in the proxy-validity check they ranked second of eight on both the proxy and the 40-epoch runs (Table~\ref{tab:proxy}). The best settings of the two searches exceeded them on the proxy by only 0.019 (PSO, 0.4953) and 0.023 (random search, 0.4995), and the four best settings of the check lie within 0.026 of each other after 40 epochs. This suggests a broad, flat region of good settings around the defaults.

Within that region, the proxy could not pick a winner. It separated poor settings from good ones ($\rho = 0.976$), but the gaps among its best candidates did not survive full training: random search's setting, the \genone proxy winner, gave the lowest \genone validation score after 100 epochs (0.643, against 0.657 for the defaults and PSO). A 12-epoch run on 40\% of the data may favor settings that learn quickly, and picking the maximum of 32 single-run scores favors scores that were high by chance; we measured neither effect.

The budget was also small: 32 evaluations in six dimensions, of which PSO spends eight on its random initial swarm and the rest in three update steps. Bergstra and Bengio~\cite{bergstra2012random} report that for most data sets only a few hyperparameters really matter, and that random search finds models as good as or better than grid search at the same budget. If only one or two of our six dimensions matter, 32 independent samples already cover them densely, leaving a guided search little to exploit. All four search winners chose an initial learning rate below the default (0.0025 to 0.0070), which fits this view, although we did not measure the importance of each dimension.

The swarm behaved as designed but converged early. In \genone, the standard deviation of its proxy scores within a round fell from 0.122 to 0.011, and its mean over all 32 evaluations exceeded that of random search (0.437 against 0.333), yet its best setting was not better. In \gentwo, PSO's best setting was the last particle of the random initial swarm, and rounds 2 to 4 never beat it (Fig.~\ref{fig:convergence}), so its chosen setting was in effect the best of eight random samples, competing with the best of 32. A global-best swarm with three update steps concentrates its later evaluations around the best point of round~1; this raises the average evaluation but not the maximum, and only the maximum is used.

Why tuning helped in \gentwo but not in \genone is less clear. Our hypothesis is that the new data, real-world photographs at 640~px with a seventh class, moved the region of good settings away from the defaults, so that any reasonable search gained about 0.02 \mapfifty on the mixed and studio test sets. It cannot be tested with the present runs, because \gentwo has no proxy run of the defaults. An alternative is training noise: with one training seed per arm, the intervals cover test sampling only (Section~\ref{subsec:limitations}), so part of a 0.02 gain could come from the particular training run. Repeating the proxy check and the full trainings for \gentwo would separate the two explanations (Section~\ref{subsec:fw-opt}).

\subsection{Implications for Tuning Claims in Waste AI}
\label{subsec:implications}

Gap~3 (Section~\ref{subsec:gaps}) states that metaheuristic tuning in this area is compared with other metaheuristics or with nothing, and that test results are reported as single numbers. Our results show what such designs can miss. Had \gentwo compared PSO only with the defaults, it would have reported a significant gain of $+0.020$ [$+0.011$, $+0.032$] \mapfifty and credited it to the swarm. Random search, at the same 32 evaluations, gained $+0.017$ [$+0.006$, $+0.030$], and the two tuned arms cannot be told apart, so the gain belongs to searching at all, not to the swarm's guidance. Single numbers would also have misled: the \genone point estimates rank the defaults first and PSO last (0.676, 0.672 and 0.668), and \gentwo validation ranked random search ahead of PSO, yet neither ordering is significant on test. Conversely, PSO's significant $+0.003$ over random search on India shows that an interval excluding zero is not enough: on scores near 0.02, such a difference has no practical meaning.

We therefore recommend a minimum protocol for any claim that a metaheuristic tunes a waste model better. First, include the defaults and random search as arms at the same number of evaluations, and report the time spent, since equal evaluations need not mean equal time (Section~\ref{subsec:res-cost}). Second, if a cheap proxy is used, check before the search that it ranks settings like longer training, against a threshold fixed in advance. Third, choose each arm's final setting on validation data and score the test split once. Fourth, report paired confidence intervals on the differences between arms, per test domain, and judge them against a difference that would matter in use. Fifth, repeat trainings and searches with several seeds. SwarmSort meets the first four points but not the fifth. In an equal-budget comparison such as that of Sayed et al.~\cite{sayed2024intelligent}, a random-search arm would cost no more than one further competitor.

\subsection{Real-World and Indian Deployment}
\label{subsec:deployment}

For a smart bin or a municipal service, the domain matters far more than the optimizer. For every arm, the real-world score is roughly half the studio score (0.298--0.310 against 0.571--0.595), and on the Indian photographs all arms collapse to 0.023--0.027 \mapfifty (Section~\ref{subsec:res-domain}). The largest tuning effect we measured, about 0.02, is more than ten times smaller than the studio-to-real-world drop of about 0.28, and it did not appear on real-world photographs at all. Local training data will therefore help more than a better search, which is why Indian data lead the roadmap (Section~\ref{subsec:fw-data}). Whether the real-world training data of \gentwo improved real-world accuracy over \genone is answered by a separate comparison on the same images: \textbf{[PENDING: Gen~1 vs Gen~2 comparison]}.

The India score must be read with its labels in mind. DWSD's boxes come from mask regions in which touching items of one class merge, and its class mapping was inferred from the masks (Section~\ref{subsec:datasets}). A model that correctly finds each bag in a pile overlaps the merged box poorly and is counted wrong, so the score measures the difference in labeling style as well as the model, in unknown proportions. A test set must be labeled like the training data before its score can be read as accuracy; relabeling part of DWSD item by item would separate the two effects.

Deployment in India also constrains the output classes. Rule~15(h) of the Solid Waste Management Rules, 2016~\cite{moefcc2016solid} requires green bins for bio-degradable waste, white bins for recyclable waste and black bins for other waste. BIODEGRADABLE maps to the green bin and the five material classes, insofar as their items are recyclable, to the white bin, but the black bin corresponds to OTHER, the weakest class of every arm (AP@50 0.050 to 0.063). A bin built on these models would be least reliable for exactly the items that belong in the black bin.

\subsection{Limitations and Threats to Validity}
\label{subsec:limitations}

Each arm was trained once (seed 42) and each search was run once (PSO seed 42, random-search seed 123), so the confidence intervals cover the sampling of test images only, not the randomness of training or search, and another seed could change the small differences between arms. No $p$-values or multiple-comparison corrections were computed for the twelve pairwise comparisons per metric in \gentwo (three pairs on four test sets), and per-class AP has no interval. The real-world test set holds a single BIODEGRADABLE box, which makes its intervals wide and lopsided.

The proxy check used eight settings, was run for \genone only, and used a 40-epoch reference rather than the 100 epochs of the final trainings; proxy validity in \gentwo, with new data, classes and image size, is assumed. The proxy also scored last-epoch weights, whereas the final models were selected as the best validation checkpoint (\texttt{best.pt}). The \gentwo Arm~C training was resumed after epoch 55 when a session froze; it continued the same weights and optimizer state but is not bit-identical to an uninterrupted run.

For DWSD, the class mapping was inferred from the masks, with six of the fifteen mask values uncertain, and the boxes come from connected regions rather than single items. Only six training hyperparameters were searched, and the budget was equal in evaluations, not in time. No external baseline was run: we did not reproduce the pipeline of Majchrowska et al.~\cite{majchrowska2022deep} or any other published detector on our data, so our scores cannot be ranked against prior work. Finally, the literature review was targeted rather than systematic (Section~\ref{subsec:review-method}), so the gaps of Section~\ref{subsec:gaps} hold for the literature we reviewed.
